\section{Completion of the fixed-length theorem}\label{sec:fixed-length}
We now collect the consequences of the full-box estimates. This also identifies exactly where their unnormalized statements become probability estimates at a specified length.

\subsection{The probability transfer}
Let \(Z_n=c_n\rho^n\) be the critical partition function of the
\(n\)-edge paths from the fixed root.
\begin{lemma}[Normalization at every length]\label{09_fast_paths:normalization}
For every integer \(n\ge0\),
\begin{equation}\label{eq:normalization}
 Z_n\ge1.
\end{equation}
\end{lemma}
\begin{proof}
The honeycomb lattice is vertex-transitive. Cutting a walk after \(m\)
edges and forgetting the mutual avoidance of its two pieces gives
\(c_{m+n}\le c_mc_n\). If \(Z_{n_0}=a<1\) for some positive
\(n_0\), writing \(n=qn_0+r\), \(0\le r<n_0\), gives
\[
 Z_n\le a^q\max_{0\le r<n_0}Z_r\le Ce^{-cn}.
\]
Trimming the two endpoint half-edges of a port bridge gives a vertex
walk from one of finitely many possible neighbors of the source port;
the weight ratio and inverse multiplicity are bounded. A height-\(h\)
bridge has at least a constant multiple of \(h\) edges. The preceding
bound would therefore imply
\(B_h\le C\sum_{n\ge c'h}e^{-cn}\le C'e^{-c''h}\), contrary to
\(B_h\asymp h^{-1/4}\). The zero-length case is immediate.
\end{proof}

Every $n$-step path from $o$ lies in a Euclidean ball of radius a constant times $n$. Choose the full-box scale $H=Kn$, with a fixed $K$ large enough to contain that ball and the bounded port extensions used in cutting. The distinction between the critical edge and vertex weights at those extensions costs a bounded factor. If a bad collection in this full box has unnormalized critical mass at most $C_AH^{-A}$, then its subset of length $n$ has $\Prob_n$-probability at most
\[
 \frac{C_AH^{-A}}{Z_n}\le C_AH^{-A}.
\]
Thus arbitrary polynomial decay in $H$ is preserved at \emph{every} length. There is no averaging over lengths in this step.

For convenience we record the three simultaneous consequences used below. For every $\eta>0$ and $k>0$, outside an event of probability $O_{\eta,k}(n^{-k})$, all subwalks $\gamma[i,j]$ and all visited centers satisfy
\begin{align}
 j-i&\le C n^\eta(1+\diam\gamma[i,j])^{4/3},\label{eq:reduction-slow}\\
 \diam\gamma[i,j]&\le C n^\eta(1+j-i)^{3/4},\label{eq:reduction-fast}\\
 M_\gamma(z,s)&\le C n^\eta s^{4/3}\quad(1\le s\le n).
 \label{eq:reduction-mass}
\end{align}
Here \eqref{eq:reduction-slow} follows from the slow-path estimate \eqref{eq:L-37}, since a directional span is bounded by the Euclidean diameter. Equation~\eqref{eq:reduction-fast} is Proposition~\ref{prop:11-modulus}, deduced from \eqref{eq:F-33b}. Equation~\eqref{eq:reduction-mass} is \eqref{eq:L-42}. The constants $C$ and the fixed factor $K^\eta$ can be absorbed by a slightly larger exponent. Those statements already count all needed positions, directions, and scales in a full box; no independence of subwalks under $\Prob_n$ is asserted.

\subsection{Diameter and local mass}
\begin{proof}[Proof of Theorem~\ref{thm:main}]
Fix $\delta,k>0$. Use \eqref{eq:reduction-slow}--\eqref{eq:reduction-mass} with a positive $\eta$ much smaller than $\delta$ and with the prescribed failure exponent $k$. All arguments below take place on their single simultaneous good event.

Applying \eqref{eq:reduction-slow} to the entire path gives
\[
 1+D\ge c n^{(3/4)(1-\eta)}.
\]
Equation~\eqref{eq:reduction-fast}, again for the entire path, gives
$D\le Cn^{3/4+\eta}$. Taking $\eta$ small enough, then $n$ large enough, absorbs the constants and the additive one and proves \eqref{eq:main-diameter}.

For the upper local mass bound, use both the deterministic bound $M_\gamma(z,s)\le n+1$ and \eqref{eq:reduction-mass}. For sufficiently small $\eta$ and large $n$ this gives
\[
 M_\gamma(z,s)\le n^\delta\min\{n,s^{4/3}\}.
\]
To prove the lower bound, fix $z=\gamma_i$ and $s\in[1,n]$, and put
$m_0=\min\{n,s^{4/3}\}$. If $n^{-\delta}m_0\le1$, the vertex $z$ itself suffices. Otherwise choose
\[
 q=\left\lfloor n^{-\delta/2}m_0\right\rfloor.
\]
For large $n$, one side of $i$ has at least $n/2\ge q$ edges available. Choose the contiguous block of $q$ edges on that side. Since $m_0>n^\delta$, we have $q\ge\tfrac12n^{-\delta/2}m_0$ and $q+1\le2n^{-\delta/2}m_0$ for large $n$. Its diameter is at most
\[
 C n^\eta(1+q)^{3/4}
 \le C' n^{\eta-3\delta/8}m_0^{3/4}
 \le C' n^{\eta-3\delta/8}s\le s,
\]
provided $\eta<3\delta/8$ and $n$ is large. The entire block lies in $\overline B(z,s)$. Self-avoidance makes its $q+1$ vertices distinct, and
$q+1\ge n^{-\delta}m_0$ for large $n$. This proves the lower local mass bound simultaneously in $i$ and $s$, using only the already simultaneous modulus.

\subsection*{Covering numbers}
For a lower bound, consider any cover by $b$ balls of radius $s$. Discard balls missing $V(\gamma)$. In each remaining ball choose one visited vertex $z$. The visited vertices in that ball lie in $\overline B(z,2s)$. When $2s\le n$, the upper local mass estimate at $2s$, with smaller exponent slack, bounds their number by
\[
 Cn^\eta\min\{n,s^{4/3}\}.
\]
When $2s>n$, the deterministic mass bound $n+1$ gives the same conclusion after changing $C$. Summing over the covering balls gives
\[
 b\ge c n^{-\eta}\max\{1,ns^{-4/3}\}.
\]
Since $\max\{1,u\}\ge(1+u)/2$, constants and the smaller exponent slack imply the lower bound in \eqref{eq:main-cover}.

For the upper bound, fix a positive $\epsilon<\delta$ and put
$m_0=\min\{n,s^{4/3}\}$ as above. If $n^{-\epsilon}m_0<2$, use one ball for each vertex. The inequality $m_0<2n^\epsilon$ implies
\[
 n+1\le Cn^\epsilon(1+ns^{-4/3}).
\]
Otherwise let $q=\lfloor n^{-\epsilon}m_0\rfloor$. Split time into consecutive blocks of at most $q$ edges, allowing consecutive blocks to share one endpoint. Each block has diameter at most
\[
 C n^\eta(1+q)^{3/4}
 \le C' n^{\eta-3\epsilon/4}s\le s
\]
if $\eta<3\epsilon/4$ and $n$ is large. A ball of radius $s$ centered at a vertex in the block covers it. The number of blocks is at most
\[
 1+\frac n q\le C n^\epsilon(1+ns^{-4/3}).
\]
Taking $\epsilon<\delta$ absorbs the constant. The whole argument is deterministic on the good event, so the bounds hold simultaneously for all real radii in the stated range. This completes the proof.
\end{proof}

Figure~\ref{fig:covering} illustrates the two deterministic steps
that turn temporal and local-mass control into covering bounds.

\begin{figure}[htbp]
\centering
\begin{tikzpicture}[scale=.8,font=\small,>=Latex]
\draw[gray!70,dashed] (0,0) circle (1.6);
\draw[blue!65!black,very thick] plot[smooth] coordinates {(-1.1,-.7) (-.7,-.5) (-.9,0) (-.3,.3) (0,0) (.5,.3) (.2,.8) (.8,.9) (1.1,.5)};
\fill (0,0) circle (2pt) node[below] {$z$};
\draw[->] (0,0)--(1.45,-.65) node[midway,below] {$s$};
\node[align=center] at (0,-2.35) {A time block of $q$ edges\\fits in one spatial ball};
\begin{scope}[xshift=6cm]
\draw[gray!70,dashed] (0,0) circle (1.3);
\draw[gray!70,dashed] (.8,.35) circle (2.6);
\fill (.8,.35) circle (2pt) node[right,xshift=12pt] {$z\in V$};
\draw[blue!65!black,very thick] plot[smooth] coordinates {(-1.1,-1.5) (-.5,-.8) (-.6,0) (.1,.1) (.8,.35) (.6,.9) (1.1,1.4)};
\node at (.4,-.65) {$B(x,s)$};
\node at (2.25,1.55) {$B(z,2s)$};
\node[align=center] at (.7,-3) {A covering ball that meets $V$\\can be recentered at a visited vertex};
\end{scope}
\end{tikzpicture}
\caption{The two deterministic uses of local control. The modulus bounds the diameter of a contiguous block and yields the upper covering bound. The local mass upper bound, after recentering, yields the lower covering bound. The curves are schematic.}
\label{fig:covering}
\end{figure}

\begin{proof}[Proof of Corollary~\ref{cor:directional-span}]
The local-length theorem holds simultaneously in the finitely many
lattice normal directions. Apply it to the whole walk after the same
exact-length probability transfer. If \(s\) is any such band span,
\(n\le Cn^\eta(1+s)^{4/3}\), so
\(s\ge c n^{(3/4)(1-\eta)}-1\). Convert band height to Euclidean
projection using \(d_0\), then choose \(\eta\) sufficiently small
relative to \(\delta\). Constants and the additive one are absorbed
for all sufficiently large \(n\).
\end{proof}

\subsection{The dimension ratio and moments}
\begin{proof}[Proof of Corollary~\ref{cor:dimension}]
Fix $\epsilon\in(0,3/4)$. On the event of Theorem~\ref{thm:main} with error $\delta<\epsilon/2$, uniformly for $s\le n^{3/4-\epsilon}$,
\[
 \log(D/s)\ge(\epsilon-\delta)\log n,
 \qquad
 \log\mathcal N(s)=\log n-\frac43\log s+O(\delta\log n)+O(1).
\]
The second assertion follows because $ns^{-4/3}\ge n^{4\epsilon/3}$; replacing $1+ns^{-4/3}$ by $ns^{-4/3}$ changes its logarithm by $O(1)$. Also
\[
 \frac43\log(D/s)=\log n-\frac43\log s+O(\delta\log n).
\]
The difference of the two quantities, divided by $\log(D/s)$, is $O(\delta/\epsilon)+o(1)$. Given a desired error, first choose $\delta$ small enough and then let $n\to\infty$. The bad event has probability tending to zero, proving the Corollary.
\end{proof}

\begin{proof}[Proof of Corollary~\ref{cor:gyration}]
The exact identity
\[
 R_g^2=\frac1{2(n+1)^2}\sum_{i,j=0}^n|\gamma_i-\gamma_j|^2
\]
shows $R_g\le D$. To get the lower bound, fix $0<\epsilon<3/4$ and put $r=n^{3/4-\epsilon}$. Use the upper local mass bound with slack $\eta<4\epsilon/3$. Uniformly in $i$, the number of $j$ within distance $r$ of $\gamma_i$ is at most
\[
 n^\eta r^{4/3}=n^{1+\eta-4\epsilon/3}=o(n).
\]
For large $n$, at least $(n+1)/2$ indices $j$ are farther than $r$ from each $\gamma_i$. Hence $R_g^2\ge r^2/4$. Arbitrarily small $\epsilon$, together with the diameter upper bound and absorption of the constant, proves the claimed probability law.

Finally both $D$ and $R_g$ are deterministically at most a constant times $n$. For fixed $p>0$, choose the failure exponent $k>p+1$. Their $p$th moments on the bad event are $O(n^{p-k})$. On the good event, their $p$th powers lie between $n^{3p/4-p\delta}$ and $n^{3p/4+p\delta}$ after adjusting the exponent error. Taking expectations and then letting $\delta\downarrow0$ proves both moment assertions.
\end{proof}
